\cotefile{136}
\scanpage{48}
\margin{\uline{\unsure{Cohomologie} des ens. $\frac12$ \unsure{simpliciaux}
et \unsure{résolutions can.} $\frac12$ \unsure{simpliciales}}}

Soit $k$ un \deleted{groupe abélien}\inserted{\unsure{anneau}}, $M$ un $k$-module,
\[
C^*:\widehat{\Delta}^{\circ}\longrightarrow
\text{\unsure{complexes croissants} de $k$-modules}
\]
un foncteur \unsure{contravariant} tel que

a) $\forall X_*\in\operatorname{Ob}\widehat{\Delta}$,
$\unsure{H^*(X_*^\vee,M)}\simeq H^*(C^*(X_*))$.
(\unsure{isom. fonctoriel} en $M$ --- \unsure{au minimum} pour $H^0$)

b) $\forall n\geq 0$, $X_*\longmapsto C^n(X_*)$
\unsure{transforme} $\unsure{\underset{\longrightarrow}{\lim}}$ en
$\unsure{\underset{\longleftarrow}{\lim}}$.

Alors b) signifie que les $C^n$ sont \unsure{représentables}
\[
C^n(X_*)=\operatorname{Hom}(X_*,C_*^n).
\]
a) implique, en \unsure{faisant} $X_*=\unsure{\Delta_n}$, que les
$C_n^*$ sont des \uline{résolutions} de $M$, \unsure{alors que l'augmentation}
$M\longrightarrow C_n^*$ est \unsure{fonctorielle} en $\unsure{\Delta_n|M}$.
$C_*^*$ est une résolution de
$\unsure{M_*=\mathbf Z_*\otimes M\simeq M_{\Delta_*}\otimes M}$,
ce qui implique bien que \unsure{l’on a un morphisme canonique}
$\unsure{H^*(C^*(X_*))\simeq H^*(X_*,M)}$,
\illegible{4} \unsure{fonctoriel}.
\margin{\unsure{Ceci est déjà} \illegible{2}!}

En prenant $X_*=\unsure{\Delta_n}$ \unsure{contractile}, on \illegible{3}
que $\unsure{H^i}$ \illegible{4}
($\unsure{H^0(C^*(X_*))=M}$).
\unsure{Pour} $\uline{\unsure{H^1}}$ \illegible{4} \unsure{pour}
$\unsure{C_*^*}$ \illegible{2}
\unsure{(analogue à} $\unsure{M_*}$\unsure{)}, i.e.
\illegible{2} $X_*\longrightarrow\unsure{Z_*^i}$

\diagram{\begin{tikzcd}
\unsure{C_*^{i-1}} \arrow[r] & \unsure{Z_*^i}\\
& X_* \arrow[u]
\end{tikzcd}}
(\unsure{pour}
\diagram{\begin{tikzcd}
C_*^0 \arrow[r] & Z_*^1\\
& X_* \arrow[u]
\end{tikzcd}})

\unsure{Si l'on} \illegible{1} $\unsure{C_*^{i-1}}$ \illegible{2},
\unsure{on n'a} \illegible{3}
\unsure{a) que les prolongements} \illegible{1} \unsure{pour les}
$C_*^i$ \unsure{sont flasques},
(\unsure{ce qui est} \illegible{1} \unsure{pour} \illegible{3}).

\deleted{\unsure{À supposer} \illegible{4} \unsure{qu'on} \illegible{2}
\unsure{pour} \illegible{1},
\unsure{fonctoriel} \illegible{4} $\unsure{C_*^i}$. On aura
\[
\unsure{H^0(X_*,C_*^0)}
\unsure{\longrightarrow}
\unsure{H^0(X_*,Z_*^1)}
\longrightarrow \unsure{H^1(X_*,M)}
\longrightarrow \unsure{H^1(X_*,C_*^0)}
\longrightarrow \unsure{H^1(X_*,Z_*^1)}.
\]
\inserted{\unsure{épim.} a)}\inserted{\unsure{inj.}}}

Mais c'est \unsure{faux déjà pour} $M=0$.
\unsure{Ceci suffit}, \unsure{pour que}, \illegible{1} b),
\unsure{que le} \illegible{2} \unsure{en} $\unsure{C_*^0\longrightarrow C_*^1}$,
\unsure{splittings} \illegible{2}!

\uline{NB} \unsure{Supposons} $C_*^*$ une \unsure{résolution flasque de} $M_*$,
les $\unsure{Z_*^{n+1}}$ \illegible{2} \unsure{connexes},
des $\unsure{K_*(M,n)}$ \illegible{1}.
\margin{47}
